\section{Algorithm: Subspace-Aware Whitening via Standard QR}

\textbf{Input:}
\begin{itemize}
    \item $G \in \mathbb{R}^{d \times B}$: Input batch (gradients or activations), with feature dimension $d$ and batch size $B$.
    \item $k$: Rank of the spike subspace (typically $k \approx d/100$).
    \item $\lambda$: Whitening strength scalar (default $1.0$).
\end{itemize}

\textbf{Output:}
\begin{itemize}
    \item $G_{clean} \in \mathbb{R}^{d \times B}$: The whitened matrix with mean bias and anisotropic spikes removed.
\end{itemize}

\subsection*{Phase 1: Mean Removal (Centering)}
Remove the Rank-1 common mean shared across all tokens.

\begin{align}
    \mu &= \frac{1}{B} G \mathbf{1}_B & \text{(Compute mean vector)} \\
    G_c &= G - \mu \mathbf{1}_B^\top & \text{(Broadcast subtract mean)}
\end{align}

\subsection*{Phase 2: Randomized Spike Detection (Sketching)}
Use a random probe matrix to identify the dominant $k$-dimensional subspace (the frequency spike).

\begin{align}
    \Omega &\sim \mathcal{N}(0, 1)^{B \times k} & \text{(Generate Gaussian random probes)} \\
    Y &= G_c \Omega & \text{(Sketch the range via Matrix Multiply)}
\end{align}
\textit{Note: $Y \in \mathbb{R}^{d \times k}$ captures the dominant directions.}

\subsection*{Phase 3: Orthogonalization via Thin QR}
Compute an orthonormal basis $Q$ for the range of $Y$ using standard thin QR decomposition. This provides better numerical stability than Cholesky QR for ill-conditioned matrices.

\begin{align}
    [Q, R] &= \text{QR}(Y, \text{mode='economic'})
\end{align}
\textit{Note: $Q \in \mathbb{R}^{d \times k}$ is the orthonormal basis for the spike. $R$ is ignored.}

\subsection*{Phase 4: Whitening (Projection \& Subtraction)}
Project the centered data onto the discovered spike basis $Q$ and subtract it to whiten the spectrum.

\begin{align}
    C &= Q^\top G_c & \text{(Compute projection coefficients)} \\
    G_{clean} &= G_c - \lambda (Q C) & \text{(Subtract reconstruction of the spike)}
\end{align}